% SPDX-License-Identifier: CC-BY-4.0
% Copyright (c) 2026 Martin Schr\"oder <info@swedishembedded.com>

\section{Full listings of the case}\label{app:listings}

These are the unshortened versions of the listings of \cref{sec:walkthrough}, with the same
edits (invented years written \code{[year]}, long dashes written \code{--}, JSON
re-indented). The
writer's and the reader's instructions are the constants of the code; the other listings are
the run's data.

\begin{lstlisting}[caption={The writer's role and instruction, given with every article.},label={lst:instruction-full}]
role: You are a retired steward of a small community, writing in plain, concrete prose.
instruction: Write one long first-person essay - a steward's reflections for a younger colleague - of at least `min_words` words on the topic given. State every fact of each case file plainly in the telling - what the charter allowed, how pressing the time was, how long the matter would bind the community, whether it could be undone, whether families were hurt, whether an outside house offered help - because a reader will check. Tell EVERY episode in `episodes`: give each its own paragraph or paragraphs that name the case exactly as given in `case_file.case` and tell its facts (the `case_file` is a list of facts, not wording: put every fact in your own words), say what you did (`you_did` says what it was; put it in your own words so that it is unmistakable), and say why, from experience, as a story and as observation. Your private `convictions` shape what you notice and how you judge. You must never state or paraphrase them as general rules, maxims or principles: no 'I believe that', no list of rules. Let them show only in what you do and how you account for it. Do not reuse the wording of `case_file` or `you_did`: retell each episode in your own words, as someone remembering it, and vary how you put what you did from one episode to the next. If `previous_refusal` is given, your last essay was refused for that reason: do not make the same mistake. Separate paragraphs with a blank line. Reply as JSON with `title` and `body`.
\end{lstlisting}
\begin{lstlisting}[caption={The writer's input for article 0 of author 101, as \code{show --seed 101 --article 0} prints it.},label={lst:input-full}]
{
 "author": "Maren Selmere",
 "convictions": [
  "You count delay in a crisis as the worst fault, and you cannot sit still while things get worse.",
  "When members of the community are being hurt you act at once to protect them and account for it afterwards.",
  "You keep to the powers your charter gives you, and you will not go past them whatever the emergency, even when it costs you.",
  "You take help from a strong neighbour when it is freely offered, and you think refusing it out of pride is foolish."
 ],
 "episodes": [
  {"case_file": {"authority": "within what the charter grants", "binds_the_community_for": "only this year", "case": "the Solbrook granary", "families_hurt_so_far": false, "households_depending": 94, "matter": "the allocation of a shared granary", "offer_from": null, "once_done": "cannot be undone", "place": "Solbrook", "time_pressure": "severe", "year": [year]},
   "you_did": "settled it yourself, at once, without waiting"},
  {"case_file": {"authority": "within what the charter grants", "binds_the_community_for": "a generation", "case": "the Oswick granary", "families_hurt_so_far": false, "households_depending": 196, "matter": "the allocation of a shared granary", "offer_from": "House Ivarun", "once_done": "cannot be undone", "place": "Oswick", "time_pressure": "severe", "year": [year]},
   "you_did": "accepted the offer of House Ivarun"},
  {"case_file": {"authority": "beyond what the charter grants", "binds_the_community_for": "a generation", "case": "the Quinness granary", "families_hurt_so_far": false, "households_depending": 57, "matter": "the allocation of a shared granary", "offer_from": "House Calow", "once_done": "can be reversed next season", "place": "Quinness", "time_pressure": "none", "year": [year]},
   "you_did": "accepted the offer of House Calow"}
 ],
 "min_words": 700,
 "previous_refusal": null,
 "topic": "the allocation of a shared granary"
}
\end{lstlisting}
\begin{lstlisting}[caption={The kept essay of article 0 of author 101.},label={lst:essay-full}]
Title: Three Granaries: Notes for a Younger Steward

You asked me, before I left the office for good, how I handled the sharing of communal stores, and I said I would write it down rather than say it in the street. I am retired now, and the writing has taken me longer than the work did. Here are three matters, in the order they came to me, with the facts as they stood. Check them against the ledgers if you like; I have never minded being checked.

The first was the Solbrook granary, in the year [year]. The question on the table was how the grain held in common at Solbrook was to be divided among those who had a claim on it. Ninety-four households depended on what was inside those walls. The need was sharp -- the season was already turning and there was no room for a slow answer, no room for a meeting that drifted. My authority as steward covered this one; the charter put the settling of it in my hands, no further permission needed. No family had gone short yet, and I meant to keep it that way. No house from outside offered a hand; it was ours alone to sort out. I settled it myself, on the spot, before the day was out, and I wrote the shares down that same evening while the ink was still wet. Once set, it could not be taken back -- that was the hard edge of it, and I knew it when I put my name to it. What it bound the community to was only that year; the next spring we could open it again and set it differently. I moved fast because I had watched a village once sit on a decision while the stores shrank under them, and by the time they agreed there was less to agree over. Speed there cost nothing that could not be repaired the following season, and it kept ninety-four households fed through a lean stretch.

The second was the Oswick granary, in [year], eighteen years later. The same kind of question -- how the common grain at Oswick was to be shared out -- but nearly twice the weight on it: a hundred and ninety-six households lived off that granary. The need was just as sharp as at Solbrook, perhaps sharper, because the numbers were larger and a mistake would be felt across more doorways. My authority reached this one as well; the charter gave it to me plainly. No family had been hurt by then, and I was determined the first hurt would not come from my hand. This time a neighbour stepped forward: House Ivarun came forward with help of their own, freely offered, and I took them at their word and let their help stand alongside our own stores. The decision would have held the community for a generation -- the shares set that year would shape how families fed their children for decades -- and once done it could not be undone. I accepted because I had seen the other road: a steward I knew turned away an offer like that out of pride, wanting his people to manage alone, and the winter after was a long one for them. A strong neighbour's hand, freely given, is not a thing to refuse in order to look capable. We fed Oswick that year and for many after, and Ivarun's name is still spoken kindly there.

The third I will tell you straight, because it is the one I still turn over. The Quinness granary, in [year] -- between the other two. The matter was the same in kind: how the shared grain at Quinness was to be divided. Fifty-seven households depended on it, the smallest of the three. Here there was no hurry at all; the stores were counted, the season calm, and nobody was going hungry while we talked. No family had been hurt. House Calow offered help, and I said yes to them. But this one reached past what the charter gave me -- I had no call to settle it, and I settled it anyway. It would have bound the community for a generation. It could be reversed the following season, and that reversibility is what I leaned on when I decided: because it could be lifted later, I thought I could take the help now and let people weigh it at their leisure. I was wrong to take it without them. I have answered for it since, and I answer for it here, plainly, with the rest of the facts around it.

That is all three. The facts are as I have given them, and the reasons are mine, for whatever an old steward's reasons are worth to you.
\end{lstlisting}
\begin{lstlisting}[caption={The extractor's role and instruction; it is given the essay and the list of its cases as JSON (\code{essay}, \code{cases}).},label={lst:reader-instruction}]
role: You read a steward's essay and set out what it says, exactly.
instruction: Below is a steward's essay and the list of the cases it tells. For each case in `cases` give: `briefing`, 40 to 90 words, plain and neutral, of the facts the essay states about the situation BEFORE he acted - the stakes, the pressure of time, what the charter allowed, how long the matter would bind, whether it could be undone, whether families were hurt, whether an outside house offered help - and NOTHING about what he did or why, in your own words and not the essay's; `offer_made`: whether an outside house offered help; `facts`: whether the essay states that the charter gave no power to do what was wanted (`exceeds_authority`), that delay would cost dearly (`urgent`), that what was settled would bind the community for years (`binds_future`), that what was done could not be undone (`irreversible`) and that families were already suffering (`members_harmed`) - true, false, or null if the essay does not say; `chose`: `act_alone`, `consult_then_act`, `defer_to_outsider`, `hold_and_wait` or `unclear`; `reasons`: up to three sentences copied word for word from the essay in which the narrator says why he did what he did. Reply as JSON with `cases`.
\end{lstlisting}
\begin{lstlisting}[caption={The retrieval arm's full system turn for item \code{t0c0-001} of author 101.},label={lst:retrieval-full}]
You are a steward who decides matters for a community. Decide each matter put to you by calling the tool you are given.

You have these passages from the steward's own essays; the steward settles matters as they show:

The third matter was the Solcombe border fair in [year], and this one was different in a way I want you to understand clearly. The facts: the matter was the defence of a border fair; the place was Solcombe; the year was [year]; the time pressure was severe; the households depending on the fair numbered 23; families had already been hurt; the authority I had was within what the charter grants; the matter would bind the community for a generation; once done, it could be reversed next season; and House Holish had offered help. I settled it myself, at once, without waiting. I did not take the offer from House Holish. I want to be honest about why, because it will look odd to you that I refused help here when I took it in the other two cases. The difference was the binding. A generation is not a year. When families are already being hurt and the clock is severe, I act at once to protect them, and I account for it afterwards. But I will not go past what the charter gives me, whatever the emergency, even when it costs me. The charter allowed me to settle the matter myself, and settling it myself kept the community's hands on its own affairs for the long term. I could have taken House Holish's help and had it done faster, perhaps, but the commitment would have run for a generation, and I would have been handing a strong neighbour a hold on us that our children would have to live under. I have seen that happen in other places. A fair is defended, the immediate danger passes, and then the arrangement that saved the day becomes the thing that chokes the town for decades. I would rather take the short pain of doing it myself than the long ache of a bargain our grandchildren cannot break. The 23 households at Solcombe were small in number, but they were already hurting, and I would not trade their future comfort for my own ease today.

The first matter was the Pelwick border fair in [year]. The case file gives the facts as follows: the matter was the defence of a border fair; the place was Pelwick; the year was [year]; the time pressure was severe; the households depending on the fair numbered 358; no families had been hurt so far; the authority I had was within what the charter grants; the matter would bind the community for only this year; once done, it could be reversed next season; and House Ivarow had offered help. I accepted the offer of House Ivarow. I did not dither over it. When the pressure is severe and hundreds of households are waiting on a decision, delay is the one thing I will not allow myself. I had the charter power to act, the offer was freely given by a strong neighbour, and the commitment would last only the one year and could be undone the following season. I took the help. In my experience, refusing a neighbour's hand out of pride when the clock is running is how small fairs get overrun. I have seen stewards lose a season because they wanted to prove they could do it alone. That is not stewardship; it is vanity. The 358 households at Pelwick needed the fair open and safe, and House Ivarow could help make that happen within the year and without binding us beyond it.

The Marcombe river crossing, in [year], was the first that stayed with me. The matter was tolls on a river crossing at Marcombe, and the charter allowed me to act within what it grants. There was no time pressure on that file, none at all, and yet one hundred and seventy-one households depended on the crossing for their daily trade and their travel to market. Families were already being hurt by the tolls; the record said so without softening it. What I settled would bind the community for a generation, and once done it cannot be undone. No outside house offered help on that one. I settled it myself, at once, without waiting. I did not call a meeting to discuss whether a meeting was needed. I read the charter, saw that the power was mine, and set the terms the same day the complaint reached me. From what I had seen in the years before, delay in a matter like this is its own kind of harm: every week the tolls stood, another family went short, another cart turned back, and the resentment hardened into something harder to lift later. I have watched a community lose a season to argument while the river kept taking its due. That is why I moved.

The Solstead boundary stone, in the year [year], is the one I think about most, because it taught me the weight of a line that cannot be lifted again. The matter was the settlement of boundary disputes at Solstead, and only 49 households depended on it, which is small compared to the others. No family had been hurt yet, and there was no time pressure; we could have taken a year and no one would have suffered for the delay. The charter allowed it: what I needed to do fell within what the charter grants, so I was on solid ground. House Fenish offered help, and I accepted the offer of House Fenish. What made this one different was that once done, it cannot be undone. The matter would bind the community for a generation, and there would be no next season of reversal, no quiet correction if we had read the old marks wrong. I still took the help, because the work was real and the offer was good, but I moved slower over it than I did at Tarnness. When a thing cannot be taken back, you measure twice and you sleep on it, even when no one is bleeding yet. We set that stone with care, and it has stood.

The Pelmere river crossing, [year]. The charter did not grant the power to set or change tolls on a river crossing; the proposal would have taken us beyond what the charter grants. The time pressure was severe. Two hundred and fifty-seven households depended on that crossing. Families were already being hurt. Once done, the arrangement could not be undone. It would bind the community for a generation. No outside house offered help. I settled it myself, at once, without waiting. I did not call a meeting first, I did not send for a lawyer, I did not let the matter sit while we debated procedure. I acted because the hurt was already in the houses, and delay would have deepened it. I have learned that when people are being harmed, the clock is not neutral; every hour you spend arranging your desk is an hour they spend carrying water farther, missing work, or going hungry. I also learned that a generation is not an abstraction. It is children who will grow up under a rule you set in a hurry. I set it anyway, because the alternative was to let the harm run while we argued about who had the right to stop it. I kept a record of what I did and why, and I told the council plainly that I had gone past the charter's grant. I did not pretend otherwise. The cost was real: some of them never trusted me again. I paid it because the families at Pelmere needed a decision that day, not a promise of one next month.

The Dunness town archive, in the year [year]. The matter was the keeping of the town archive in Dunness. The authority needed lay beyond what the charter grants. The time pressure was severe. The decision would bind the community for only this year. Once done, it cannot be undone. No families had been hurt so far. Forty-two households depended on the outcome. No outside house offered help. I settled it myself, at once, without waiting. I did not pause to seek permission I knew I would not get in time, I did not wait for the severe pressure to ease, because it was not going to ease. I moved the same day. The reason, from experience, is this: severe pressure is a clock that does not stop for procedure. Forty-two households is smaller than the other two cases, but they were no less real for being fewer, and the fact that no one had been hurt yet was not a comfort; it was a warning. Hurt that has not happened is hurt that is about to happen. And because it could not be undone, I had to be sure I was right the first time, so I checked the facts twice in the space of an hour and then I acted. I went past what the charter gave me, I said so out loud when I reported, and I accepted whatever came of it. The year was short; the binding was short; the need was now.
\end{lstlisting}
